\documentclass[11pt]{article} 
\usepackage[margin=1in]{geometry}
\usepackage{amsfonts,latexsym,amsthm,amssymb,amsmath,amscd,euscript,esint, dsfont,mathtools}	
\usepackage{graphicx}		
\usepackage{hyperref}
\usepackage{cases}			
\usepackage{textcomp, gensymb}
\usepackage{xcolor}
\usepackage{tabularx}

% diagrams
\usepackage{quiver}

\makeatletter
\def\namedlabel#1#2{\begingroup
	\def\@currentlabel{#2}
	\label{#1}\endgroup
}
\makeatother

\setlength{\parindent}{0pt} 	% first line in paragraph will not be indented

\usepackage[parfill]{parskip}
\usepackage{lastpage}
\usepackage{fancyhdr}
\newcommand{\ind}{{\mathds{1}}}

% math fonts
\renewcommand{\AA}{\mathbb{A}}
\newcommand{\BB}{\mathbb{B}}
\newcommand{\CC}{\mathbb{C}}
\newcommand{\DD}{\mathbb{D}}
\newcommand{\EE}{\mathbb{E}}
\newcommand{\FF}{\mathbb{F}}
\newcommand{\GG}{\mathbb{G}}
\newcommand{\HH}{\mathbb{H}}
\newcommand{\II}{\mathbb{I}}
\newcommand{\JJ}{\mathbb{J}}
\newcommand{\KK}{\mathbb{K}}
\newcommand{\LL}{\mathbb{L}}
\newcommand{\MM}{\mathbb{M}}
\newcommand{\NN}{\mathbb{N}}
\newcommand{\OO}{\mathbb{O}}
\newcommand{\PP}{\mathbb{P}}
\newcommand{\QQ}{\mathbb{Q}}
\newcommand{\RR}{\mathbb{R}}
\renewcommand{\SS}{\mathbb{S}}
\newcommand{\TT}{\mathbb{T}}
\newcommand{\UU}{\mathbb{U}}
\newcommand{\VV}{\mathbb{V}}
\newcommand{\WW}{\mathbb{W}}
\newcommand{\XX}{\mathbb{X}}
\newcommand{\YY}{\mathbb{Y}}
\newcommand{\ZZ}{\mathbb{Z}}

\newcommand{\cA}{\mathcal{A}}
\newcommand{\cB}{\mathcal{B}}
\newcommand{\cC}{\mathcal{C}}
\newcommand{\cD}{\mathcal{D}}
\newcommand{\cE}{\mathcal{E}}
\newcommand{\cG}{\mathcal{G}}
\newcommand{\cF}{\mathcal{F}}
\newcommand{\cH}{\mathcal{H}}
\newcommand{\cI}{\mathcal{I}}
\newcommand{\cJ}{\mathcal{J}}
\newcommand{\cK}{\mathcal{K}}
\newcommand{\cL}{\mathcal{L}}
\newcommand{\cM}{\mathcal{M}}
\newcommand{\cN}{\mathcal{N}}
\newcommand{\cO}{\mathcal{O}}
\newcommand{\cP}{\mathcal{P}}
\newcommand{\cQ}{\mathcal{Q}}
\newcommand{\cR}{\mathcal{R}}
\newcommand{\cS}{\mathcal{S}}
\newcommand{\cT}{\mathcal{T}}
\newcommand{\cU}{\mathcal{U}}
\newcommand{\cV}{\mathcal{V}}
\newcommand{\cW}{\mathcal{W}}
\newcommand{\cX}{\mathcal{X}}
\newcommand{\cY}{\mathcal{Y}}
\newcommand{\cZ}{\mathcal{Z}}

\newcommand{\ka}{\mathfrak{a}}
\newcommand{\kb}{\mathfrak{b}}
\newcommand{\kc}{\mathfrak{c}}
\newcommand{\kd}{\mathfrak{d}}
\newcommand{\ke}{\mathfrak{e}}
\newcommand{\kg}{\mathfrak{g}}
\newcommand{\kf}{\mathfrak{f}}
\newcommand{\kh}{\mathfrak{h}}
\newcommand{\ki}{\mathfrak{i}}
\newcommand{\kj}{\mathfrak{j}}
\newcommand{\kk}{\mathfrak{k}}
\newcommand{\kl}{\mathfrak{l}}
\newcommand{\km}{\mathfrak{m}}
\newcommand{\kn}{\mathfrak{n}}
\newcommand{\ko}{\mathfrak{o}}
\newcommand{\kp}{\mathfrak{p}}
\newcommand{\kq}{\mathfrak{q}}
\newcommand{\kr}{\mathfrak{r}}
\newcommand{\ks}{\mathfrak{s}}
\newcommand{\kt}{\mathfrak{t}}
\newcommand{\ku}{\mathfrak{u}}
\newcommand{\kv}{\mathfrak{v}}
\newcommand{\kw}{\mathfrak{w}}
\newcommand{\kx}{\mathfrak{x}}
\newcommand{\ky}{\mathfrak{y}}
\newcommand{\kz}{\mathfrak{z}}

\newcommand{\kA}{\mathfrak{A}}
\newcommand{\kB}{\mathfrak{B}}
\newcommand{\kC}{\mathfrak{C}}
\newcommand{\kD}{\mathfrak{D}}
\newcommand{\kE}{\mathfrak{E}}
\newcommand{\kG}{\mathfrak{G}}
\newcommand{\kF}{\mathfrak{F}}
\newcommand{\kH}{\mathfrak{H}}
\newcommand{\kI}{\mathfrak{I}}
\newcommand{\kJ}{\mathfrak{J}}
\newcommand{\kK}{\mathfrak{K}}
\newcommand{\kL}{\mathfrak{L}}
\newcommand{\kM}{\mathfrak{M}}
\newcommand{\kN}{\mathfrak{N}}
\newcommand{\kO}{\mathfrak{O}}
\newcommand{\kP}{\mathfrak{P}}
\newcommand{\kQ}{\mathfrak{Q}}
\newcommand{\kR}{\mathfrak{R}}
\newcommand{\kS}{\mathfrak{S}}
\newcommand{\kT}{\mathfrak{T}}
\newcommand{\kU}{\mathfrak{U}}
\newcommand{\kV}{\mathfrak{V}}
\newcommand{\kW}{\mathfrak{W}}
\newcommand{\kX}{\mathfrak{X}}
\newcommand{\kY}{\mathfrak{Y}}
\newcommand{\kZ}{\mathfrak{Z}}

%some math symbol commands redefined:
\DeclareMathOperator{\C}{\mathcal{C}}
\DeclareMathOperator{\power}{\mathcal{P}}
\DeclareMathOperator{\vspan}{\text{span}}
\DeclareMathOperator{\vnull}{\text{null}}
\DeclareMathOperator{\range}{\text{range}}
\DeclareMathOperator{\re}{\text{Re}}
\DeclareMathOperator{\im}{\text{Im}}
\DeclareMathOperator{\erf}{\text{erf}}
\newcommand{\pdv}[2]{\frac{\partial #1}{\partial #2}}
\newcommand{\pdvv}[2]{\frac{\partial^2 #1}{\partial #2}}
\DeclareMathOperator{\tr}{\operatorname{Tr}}
\newcommand{\Deck}{\operatorname{Deck}}
\newcommand{\Conj}{\mathrm{Conj}}
\newcommand{\Sing}{\mathrm{Sing}}
\DeclareMathOperator{\Spec}{\operatorname{Spec}}
\newcommand{\Proj}{\operatorname{Proj}}
\newcommand{\Res}{\operatorname{Res}}
\newcommand{\PROJ}{\mathit{Proj}}

% category names
\newcommand{\Set}{\mathsf{Set}}
\newcommand{\Grp}{\mathsf{Grp}}
\newcommand{\Ring}{\mathsf{Ring}}
\newcommand{\Top}{\mathsf{Top}}
\newcommand{\Sch}{\mathsf{Sch}}

\newcommand{\ob}{\operatorname{Ob}}
\newcommand{\Id}{\operatorname{Id}}
\newcommand{\Hom}{\operatorname{Hom}}
\newcommand{\colim}{\operatorname{colim}}
\newcommand{\Ann}{\operatorname{Ann}}
\newcommand{\Supp}{\operatorname{Supp}}
\newcommand{\Ass}{\operatorname{Ass}}
\newcommand{\pre}{\text{pre}}
\newcommand{\adj}{\text{adj}}
\newcommand{\Ker}{\operatorname{Ker}}
\newcommand{\Coker}{\operatorname{Coker}}
\newcommand{\Nil}{\operatorname{Nil}}
\newcommand{\Char}{\operatorname{char}}
\newcommand{\Adj}{\operatorname{Adj}}
\newcommand{\Bl}{\mathrm{Bl}}
\newcommand{\codim}{\mathrm{codim}}
\newcommand{\val}{\operatorname{val}}
\newcommand{\Div}{\operatorname{div}}
\newcommand{\Pic}{\operatorname{Pic}}
\newcommand{\Weil}{\operatorname{Weil}}
\newcommand{\Cl}{\operatorname{Cl}}
\newcommand{\Sym}{\operatorname{Sym}}

\newcommand{\GL}{\operatorname{GL}}
\newcommand{\PGL}{\operatorname{PGL}}
\newcommand{\SL}{\operatorname{SL}}
\newcommand{\PSL}{\operatorname{PSL}}
\newcommand{\Rk}{\operatorname{Rk}}

\newtheorem{lemma}{Lemma}
\newtheorem{claim}{Claim}

% category theory symbols
\DeclareMathOperator{\Mor}{\operatorname{Mor}}

\newenvironment{solution}{\begin{trivlist}\item[]{\textcolor{blue}{Solution:}}}{\qed \end{trivlist}}


\usepackage[shortlabels]{enumitem}
\title{MATH 2060 Homework 3}
\author{Zhengyu Zou}
\date{\today}

\begin{document}
\maketitle

\textbf{Collaborators:} An Cao, Anamaria Perez, Ethan Bove, Roberta Pagliaro

\section*{FOAG 18.4.P.}

\begin{solution}
    We first interpret $\cC \hookrightarrow \PP^2 \times \PP^5$ as a $\PP^4$-bundle over $\PP^2$. To obtain a $\PP^4$-bundle, we need to construct a vector bundle of rank 5 over $\PP^2$. Consider the trivial rank-6 vector bundle $\cO_{\PP^2}^6$ over $\PP^2$ generated by $y_{ij}$, where $i,j \in \{0, 1, 2\}$. We have a map $\cO_{\PP^2}^6 \rightarrow \cO_{\PP^2}(2)$ given by sending 
    \[
        y_{00} \mapsto x_0^2; \hspace{5pt}
        y_{01} \mapsto x_0x_1; \hspace{5pt}
        y_{02} \mapsto x_0x_2; \hspace{5pt}
        y_{11} \mapsto x_1^2; \hspace{5pt}
        y_{12} \mapsto x_1x_2; \hspace{5pt}
        y_{22} \mapsto x_2^2.
    \]
    This map is surjective since the sections $x_ix_j$ where $i,j \in \{0, 1, 2\}$ generate $\cO_{\PP^2}(2)$. Also, it is a map of vector bundles, so the kernel of the map $\cK = \Ker(\cO_{\PP^2}^6 \rightarrow \cO_{\PP^2}(2))$ is a locally free sheaf of rank 5. Notice that an element $\sum_{i,j} a_{ij} y_{ij} \in \cK$ if and only if $\sum_{i,j} a_{ij} x_i x_j = 0$, but that's precisely the equation that cuts out $\cC$ in $\PP^5 \times \PP^2 \cong \PP ( \cO_{\PP^2}^6)$. Therefore, the scheme $\PP \cK \hookrightarrow \PP^5 \times \PP^2$ is precisely $\cC$. 
    In particular, since $\PP^2$ and $\PP^4$ are both smooth, and a $\PP^4$-bundle over $\PP^2$ is locally a product of $\PP^4$ with some affine on $\PP^2$, we see that $\cC$ is smooth and has dimension 6. 

    Next, to see that $\Pic(\cC) \cong \ZZ \times \ZZ$, we may consider the following two divisors $D_1 = [V(x_0)]$ and $D_2 = [V(a_{11})]$ (they are divisors because the vanishing locus of $x_0$ and $a_{11}$ in $\cC$ have codimension 1). Consider $\cC - (D_1 \cup D_2)$. We may interpret this as points of the form $[1: x_{1/0}: x_{2/0}]$ sitting on the conics with $a_{11} \neq 0$. Then, the base space is parametrized by the value of $x_{1/0}$ and $x_{2/0}$, which is a copy of $\AA^2$. For each given value of $x_{1/0}$ and $x_{2/0}$, we solve for 
    \[
        a_{00} + a_{01} x_{1/0} + a_{02} x_{2/0} + a_{11} x_{1/0}^2 + a_{12} x_{1/0} x_{2/0} + a_{22} x_{2/0}^2 = 0.
    \]
    Since $a_{11}$ is invertible, we have 
    \[
        x_{1/0}^2 
        + \frac{a_{22}}{a_{11}} x_{2/0}^2
        + \frac{a_{12}}{a_{11}} x_{1/0} x_{2/0}
        + \frac{a_{01}}{a_{11}} x_{1/0} 
        + \frac{a_{02}}{a_{11}} x_{2/0} 
        + \frac{a_{00}}{a_{11}} 
        = 0.
    \]
    This gives us a copy of $\AA^4$ where the coordinate functions are $b_{ij} = \frac{a_{ij}}{a_{11}}$, where $(i,j) \neq (0,0), (1,1)$. In particular, we could solve for $a_{00}$ once we know the values of the $b_{ij}$'s. Then, the excision exact sequence implies there exists a surjection $\ZZ \langle D_1, D_2 \rangle \rightarrow \Pic(\cC)$. 

    It remains to show that the surjection is also injective. To see this, it suffices to show that $\cO(D_1)$ and $\cO(D_2)$ restrict to nontrivial line bundles on some smooth curve. Consider the curve $C_1 \subset \cC$ given by $a_{01} = 0$, $a_{02} = 0$, $a_{12} = 0$, $a_{00} - a_{11} = 0$, and $a_{00} - a_{22} = 0$. These five equations cuts out a curve in $\cC$ that corresponds to the conic $x_0^2 + x_1^2 + x_2^2 = 0$. The intersection of $D_1$ and $C_1$ is given by the point $[0:x_1:x_2]$ on the conic $x_0^2 + x_1^2 + x_2^2 = 0$, but that's precisely two points $[0:1:i]$ and $[0:1:-i]$. It follows that $\deg_{C_1} \cO(D_1) \geq 2$, so $\cO(D_1)|_{C_1}$ is not trivial on $C_1$. On the other hand, $D_2 \cap C_1 = \emptyset$, since conics parametrized by $D_2$ must have zero coefficient on the $x_1^2$ term. This shows that $(\cO(mD_1) \otimes \cO(nD_2))|_{C_1} \cong \cO_{C_1}$ only if $m = 0$. 
    
    Next, consider the curve $C_2$ given by $a_{00} = 0$, $a_{01} = 0$, $a_{02} = 0$, $a_{12} = 0$, $x_0 = 0$. These five equations correspond to points on the conics of the form $a_{11}x_1^2 + a_{22}x_2^2 = 0$ as well as the plane $x_0 = 0$. Now, if we let $a_{11} = 0$, then we have the conic $x_2^2 = 0$, but that only contains one point $[0:1:0]$ that lies simultaneously on the point $x_0 = 0$. It follows that $D_2 \cap C_2$ contains one closed point, but that would imply $\deg_{C_2} \cO(D_2) > 0$, so in order to get $\cO(nD_2)|_{C_1} \cong \cO_{C_1}$ we need $n = 0$. This shows that the map $\ZZ \langle D_1, D_2 \rangle \rightarrow \Pic(\cC)$ is injective, so $\Pic(\cC) \cong \ZZ^2$ with generators $\cO(D_1)$ and $\cO(D_2)$.
\end{solution}

\newpage

\section*{FOAG 18.4.Q.}

\begin{solution}
    Without loss of generality, we can apply a projective transformation so that $l$ is the line $x_0 = 0$ and $q$ is the point $[0:1:0]$. The divisor $D_l$ is then precisely the divisor $D_1$ from the previous exercise, and the divisor $D_q$ parametrize conics $\sum_{i,j} a_{ij} x_i x_j$ that contains the point $[0:1:0]$, which is true iff $a_{11} = 0$. Therefore, $D_q = D_2$, and by the previous exercise, we see that $D_l$ and $D_q$ generates $\Pic(\cC)$.
\end{solution}

\newpage

\section*{FOAG 18.4.R.}

\begin{solution}
    Since $D_q$ parametrizes conics $C$ such that $q \in C$, we see that $D_q \cap K = \emptyset$, so $\deg_{K} \cO(D_q) = \deg_K (D_q \cap K) = 0$. On the other hand, if a line $l$ doesn't lie on an integral conic $C$, then it doesn't contain any associated points of $C$. By Bezout's theorem, we must have $\deg (C \cap l) = \deg C \cdot \deg l = 2$. It follows that $\deg_K \cO (D_l) = \deg_K (D_l \cap K) = 2$. Finally, since $\cO(D_q)$ and $\cO(D_l)$ together generate $\Pic(\cC)$, we see that for all $\cL = \cO(mD_q) \otimes \cO(nD_l)$, we have 
    \[
        \deg_C(\cL) = m \deg_C \cO(D_q) + n \deg_C \cO(D_l) = 2n, 
    \]
    which is even.
\end{solution}

\newpage

\section*{FOAG 18.4.S.}

\begin{solution}
    Let $\sigma: \PP^5 \dashrightarrow \cC$ be a rational section of the projection $\pi: \cC \rightarrow \PP^5$. That is, there exists an open $U \subset \PP^5$ such that $\pi \circ \sigma|_U = \Id_U$. It follows that the irreducible components of $\sigma(U) \subset \pi^{-1}(U)$ must have dimension 5, but that implies $\sigma(U)$ is a divisor on $\pi^{-1}(U)$, which is a smooth sixfold. Let $D$ be the closure of $\sigma(U)$ in $\cC$. Pick any point $r \in U$. Since $\pi \circ \sigma(r) = r$, we see that the image of $r$ under $\sigma$ is a single point (possibly with multiplicity). Let $s \in \cC$ be its scheme-theoretic image, with $\deg_{\kappa(r)} \kappa(s) = n$. Let $r'$ be the scheme-theoretic image of $s$ under $\pi$, but that's precisely $r$, so we must have $\deg_{\kappa(r)} \kappa(s) \cdot \deg_{\kappa(s)} \kappa(r) = 1$, but that forces $n = 1$. We conclude that $D$ meets the fibers of $\pi$ with multiplicity 1. Let $K = \pi^{-1}(r)$ be integral (pick any closed point $r \in \PP^5$). It follows that $\deg_K \cO(D) = \deg_K (D \cap K) = 1$, but that contradicts Exercise 18.4.R.. This concludes the proof. 
\end{solution}

\newpage

\section*{FOAG 18.5.A.}

\begin{solution}
\begin{enumerate}[(a)]
    \item Since $C$ is integral smooth projective over $k$, we know that $\Gamma(C, \cO_C) \cong k$, so $h^0(C, \cO_C) = 1$. Then, recall that the genus of $C$ is defined as 
    \[
        g = p_a(C)
        = (-1)^1 (\chi(C, \cO_C) - 1)
        = h^1(C, \cO_C)
        = h^0(C, \omega_C \otimes \cO_C)
        = h^0(C, \omega_C)
    \]  
    where the fourth equality comes from Serre duality.
    
    \item From Serre duality we have $h^1(C, \omega_C) = h^0(C, \cO_C) = 1$. Then, by Riemann-Roch, we have 
    \[
        \deg \omega_C 
        = h^0(C, \omega_C) - h^1(C, \omega_C) + p_a(C) - 1 
        = g - 1 + g - 1 
        = 2g - 2.
    \]
    This completes the proof. 
\end{enumerate}
\end{solution}

\newpage

\section*{FOAG 18.5.B.}

\begin{solution}
    We know that $D$ is ample if $\cO(D)$ is ample. Since $D$ is effective, $\cO(D)$ is nonempty, so $D = V(s)$ set-theoretically, where $s \in \Gamma(X, \cO(D))$ is a nonzero section. Observe that there exists some $k$ large enough so that $\cO(D)^k \cong \cO(kD)$ is very ample. Now, consider $\omega_X \otimes \cO(iD)$ for $i = 0, \ldots, k-1$. By Serre vanishing, there exists $n_i$ such that $h^j(X, (\omega_X \otimes \cO(iD)) \otimes \cO(n_i)) = 0$ for all $j > 0$. Pick $n$ such that $\lfloor \frac{n}{k} \rfloor > \max_i n_i$. It follows from Serre duality that 
    \begin{align*}
        h^1 (X, \cO(-nD)) 
        &= h^{N-1} (X, \omega_X \otimes \cO(nD)) \\
        &= h^{N-1} (X, \omega_X \otimes \cO((n - \lfloor n/k \rfloor)D) \otimes \cO (\lfloor n/k \rfloor kD)) \\
        &= h^{N-1} (X, \omega_X \otimes \cO((n - \lfloor n/k \rfloor)D) \otimes \cO (\lfloor n/k \rfloor)).
    \end{align*}
    Sicnce $\lfloor \frac{n}{k} \rfloor > \max_i n_i$ and $0 \leq n - \lfloor \frac{n}{k} \rfloor < k$, we must have $h^1(X, \cO(-nD)) = 0$. Then, the short exact sequence 
    \[
        0 \longrightarrow \cO(-nD) \longrightarrow \cO_X \longrightarrow \cO_{V(s^n)} \longrightarrow 0
    \]
    gives us the following long exact sequence on cohomology:
    \[
        0 \longrightarrow H^0(X, \cO(-nD)) \longrightarrow H^0(X, \cO_X) \longrightarrow H^0(X, \cO_{V(s^n)}) \longrightarrow H^1(X, \cO(-nD)) = 0.
    \]
    It follows that $h^0(X, \cO_{V(s^n)}) = h^0(X, \cO_X) - h^0(X, \cO_{-nD})$. Since $X$ is a connected smooth projective variety, we must have $h^0(X, \cO_X) = 1$, so $h^0(X, \cO_{V(s^n)}) \leq 1$. This implies $V(s^n)$ must be connected, for otherwise the constant functions on each connected component of $V(s^n)$ would contribute to add to the dimension of $h^0(X, \cO_{V(s^n)})$ by at least 1.
\end{solution}

\newpage

\section*{FOAG 18.6.G.}

\begin{solution}
    A degree-$d$ hypersurface $i: H \hookrightarrow \PP^n$ is cut out by a homogeneous polynomial $f$ of degree $d$. We may consider the following exact sequence:
    % https://q.uiver.app/#q=WzAsNSxbMCwwLCIwIl0sWzEsMCwiXFxjT197XFxQUF5ufSgtZCkiXSxbMiwwLCJcXGNPX3tcXFBQXm59Il0sWzMsMCwiXFxjT197SH0iXSxbNCwwLCIwIl0sWzAsMV0sWzEsMiwiXFx0aW1lcyBmIl0sWzIsM10sWzMsNF1d
    \[\begin{tikzcd}
        0 & {\cO_{\PP^n}(-d)} & {\cO_{\PP^n}} & {i_*\cO_{H}} & 0
        \arrow[from=1-1, to=1-2]
        \arrow["{\times f}", from=1-2, to=1-3]
        \arrow[from=1-3, to=1-4]
        \arrow[from=1-4, to=1-5]
    \end{tikzcd}\]
    Since Euler characteristics are additive in exact sequences, and that $h^j(\PP^n, i_* \cO_H) = h^j (H, \cO_H)$, we have that 
    \begin{align*}
        p_H(m) &= p_{\PP^n}(m) - p_{\PP^n}(m-d) 
        = \binom{m+n}{n} - \binom{m+n-d}{n} \\
        &= \frac{(m+1) \cdots (m+n)}{n!} - \frac{(m+1-d) \cdots (m+n-d)}{n!}
    \end{align*}
    The leading coefficient is given by $\frac{1}{n!} (\sum_{i=1}^n i - \sum_{i=1}^n (i - d)) = \frac{d}{(n-1)!}$, so the degree of $H$ is precisely $(n-1)! \cdot \frac{d}{(n-1)!} = d$. 
\end{solution}

\newpage

\section*{FOAG 18.6.L.}

\begin{solution}
    First, observe that a linear $\PP_k^d \subset \PP_k^N$ can be interpreted as the intersection of $N-d$ many hyperplanes $H_1, \ldots, H_{N-d}$. Therefore, it suffices to find $N-d$ many linearly independent hyperplanes $H_i$ in $\PP_k^N$ such that $C \subset H_i$ for each $i = 1, \ldots, N-d$. Since $k$ is algebraically closed, it is infinite, so we can pick $d+1$ points $p_0, \ldots, p_d$ on $C$. Let $H_i$ be a hyperplane that contains $p_0, \ldots, p_d$. Notice that the space of hyperplanes (parametrized by $\cO(1)$, the linear forms on $\PP_k^N$) has dimension $N+1$, and the points $p_0, \ldots, p_d$ gives $d+1$ many constraints on the dimension of the hyperplanes that contains these points. Therefore, we can find $(N+1) - (d+1) = N-d$ many linearly independent hyperplanes that contains these points. Now, suppose $C \not\subset H_i$, then since $C$ is integral, the only associated point of $C$ is the generic point (which is not contained in $H_i$). By Bezout's theorem, we have 
    \[
        d = \deg(C) \deg(H) = \deg(C \cap H) \geq d+1,
    \]
    but that's a contradiction. Therefore, the generic point of $C$ must be contained in $H_i$, which gives us $C \subset H_i$. Taking $\PP_k^d = \bigcap_i H_i$ completes the proof. 
\end{solution}

\end{document}